% ============================================================
% Section 2: Mathematical Foundation
% ============================================================
\section{Mathematical Foundation}

\begin{frame}{The Mathematics (Brief)}
    \textbf{1. Kernel Density Estimation:}
    Given $n$ data points $\{x_1, x_2, \ldots, x_n\}$, the kernel density estimate at point $x$ is:

    \[
    \hat{f}(x) = \frac{1}{n h} \sum_{i=1}^{n} K\left(\frac{x - x_i}{h}\right)
    \]

    where we use a \textbf{symmetric kernel}: $K(x) = c_k \cdot k(x^2)$
    \begin{itemize}
        \item $k: \mathbb{R}^+ \to \mathbb{R}$ is the \textbf{kernel profile} (function of squared distance)
        \item $h > 0$ is the bandwidth parameter
        \item $c_k$ is a normalization constant
        \item $x_i$ are the data points and no need to be inside the bandwidth
    \end{itemize}

    \small{\textcolor{customred}{Key insight:} If we set $h$ too large, we may suffer from over-smoothing.}
\end{frame}

\begin{frame}{From KDE to Mean Shift}
    \textbf{2. Deriving the Mean Shift Vector:}

    To find density modes, we take the \textbf{gradient} of $\hat{f}(x)$ and set it to zero.

    \vspace{0.2cm}
    Define the \textbf{shadow kernel profile}: $g(s) = -k'(s)$ (negative derivative of $k$)

    \vspace{0.2cm}
    Taking $\nabla \hat{f}(x) = 0$ and rearranging yields:
    \[
    x = \frac{\sum_{i=1}^{n} x_i \, g\left(\left\|\frac{x - x_i}{h}\right\|^2\right)}{\sum_{i=1}^{n} g\left(\left\|\frac{x - x_i}{h}\right\|^2\right)}
    \]

    \small{\textcolor{customred}{Key insight:} The function $g = -k'$ ensures we move \textbf{against} the gradient direction (toward increasing density). This weighted mean tells us where the mode is!}
\end{frame}

\begin{frame}{The Mathematics (Contd.)}

    \textbf{3. Mean Shift Vector:}

    The difference between the weighted mean and current position:
    \[
    m(x) = \underbrace{\frac{\sum_{i=1}^{n} x_i \, g\left(\left\|\frac{x - x_i}{h}\right\|^2\right)}{\sum_{i=1}^{n} g\left(\left\|\frac{x - x_i}{h}\right\|^2\right)}}_{\text{weighted mean using } g} - x
    \]

    \vspace{0.3cm}

    \textbf{4. Update Rule:}
    \[
    x^{(t+1)} = x^{(t)} + m(x^{(t)})
    \]

    \small{\textcolor{customred}{Key insight:} Since $g = -k'$, the mean shift vector always points toward increasing density (gradient ascent on the density surface).}
\end{frame}

\begin{frame}{Common Kernel Profiles}
    \small
    Let $s = \left\|\frac{x - x_i}{h}\right\|^2$ be the squared normalized distance.
    \vspace{0.2cm}
    \begin{columns}[T]
        \begin{column}{0.5\textwidth}
            \textbf{Gaussian:}
            \begin{align*}
            k(s) &= \exp\left(-\frac{s}{2}\right) \\
            g(s) &= -k'(s) = \frac{1}{2}\exp\left(-\frac{s}{2}\right)
            \end{align*}

            \vspace{0.3cm}

            \textbf{Epanechnikov:}
            \begin{align*}
            k(s) &= \begin{cases} 1 - s & s \leq 1 \\ 0 & \text{else} \end{cases} \\
            g(s) &= \begin{cases} 1 & s \leq 1 \\ 0 & \text{else} \end{cases}
            \end{align*}
        \end{column}

        \begin{column}{0.5\textwidth}
            \textbf{Flat (Uniform):}
            \begin{align*}
            k(s) &= \begin{cases} 1 & s \leq 1 \\ 0 & \text{else} \end{cases} \\
            g(s) &= 0 \text{ (undefined at boundary)}
            \end{align*}

            \vspace{0.3cm}

            \small{\textcolor{customred}{Why Epanechnikov?}}
            \begin{itemize}
                \item $g(s) = 1$ (constant) $\Rightarrow$ uniform weights
                \item Simple weighted mean of neighbors
                \item Most efficient computation
                \item Used by sklearn's MeanShift
            \end{itemize}

            \vspace{0.2cm}
            \small{\textcolor{customred}{Why not Flat?}}
            \begin{itemize}
                \item $g(s) = 0$ everywhere $\Rightarrow$ no gradient info!
            \end{itemize}
        \end{column}
    \end{columns}
\end{frame}


\begin{frame}{Algorithm Steps}
    \begin{algorithm}[H]
        \caption{Mean Shift Clustering}
        \begin{algorithmic}[1]
            \State \textbf{Input:} Data points $\{x_1, \ldots, x_n\}$, bandwidth $h$, kernel profile $k$, $\epsilon$ for convergence
            \State \textbf{Output:} Cluster assignments
            \State
            \For{each data point $\mathbf{x}_i$}
                \State $\mathbf{y} \gets \mathbf{x}_i$ \Comment{Initialize at data point}
                \While{not converged}
                    \State Compute mean shift: $\mathbf{m}(\mathbf{y})$
                    \State Update: $\mathbf{y} \gets \mathbf{y} + \mathbf{m}(\mathbf{y})$
                \EndWhile
                \State Store converged point (mode) for $\mathbf{x}_i$
            \EndFor
            \State
            \State Group points by their converged modes
            \State \textbf{Return} cluster assignments
        \end{algorithmic}
    \end{algorithm}
\end{frame}
